%; whizzy paragraph -pdf xpdf -latex ./whizzypdfptex.sh
%; whizzy-paragraph "^\\\\begin{frame}"
% latex beamer presentation.
% platex, latex-beamer でコンパイルすることを想定。 

%     Tokyo Debian Meeting resources
%     Copyright (C) 2009 Junichi Uekawa
%     Copyright (C) 2009 Nobuhiro Iwamatsu

%     This program is free software; you can redistribute it and/or modify
%     it under the terms of the GNU General Public License as published by
%     the Free Software Foundation; either version 2 of the License, or
%     (at your option) any later version.

%     This program is distributed in the hope that it will be useful,
%     but WITHOUT ANY WARRANTY; without even the implied warreanty of
%     MERCHANTABILITY or FITNESS FOR A PARTICULAR PURPOSE.  See the
%     GNU General Public License for more details.

%     You should have received a copy of the GNU General Public License
%     along with this program; if not, write to the Free Software
%     Foundation, Inc., 51 Franklin St, Fifth Floor, Boston, MA  02110-1301 USA

\documentclass[cjk,dvipdfmx,12pt]{beamer}
\usetheme{Tokyo}
\usepackage{monthlypresentation}

%  preview (shell-command (concat "evince " (replace-regexp-in-string "tex$" "pdf"(buffer-file-name)) "&")) 
%  presentation (shell-command (concat "xpdf -fullscreen " (replace-regexp-in-string "tex$" "pdf"(buffer-file-name)) "&"))
%  presentation (shell-command (concat "evince " (replace-regexp-in-string "tex$" "pdf"(buffer-file-name)) "&"))

%http://www.naney.org/diki/dk/hyperref.html
%日本語EUC系環境の時
\AtBeginDvi{\special{pdf:tounicode EUC-UCS2}}
%シフトJIS系環境の時
%\AtBeginDvi{\special{pdf:tounicode 90ms-RKSJ-UCS2}}

\title{DebianでSphinxとDoxygenを使ってみた}
%\subtitle{第75回 2011年6月度}
\author{まえだこうへい mkouhei@debian.or.jp\\IRC nick: mkouhei}
\date{2011年6月18日}
\logo{\includegraphics[width=8cm]{image200607/openlogo-light.eps}}

\begin{document}

\frame{\titlepage{}}

\section{}

\begin{frame}
 \frametitle{最近の流行りらしい}
  \begin{itemize}
   \item \url{http://sphinx-users.jp/example.html}
  \end{itemize}
\end{frame}

\begin{frame}{reSTとSphinxの概要}
 \begin{itemize}
  \item reST(reStructuredText)から各種フォーマットのドキュメントを生成
  \item HTML、 \LaTeX 、PDF、ePub、man、平文テキスト、JSONなど
 \end{itemize}
\end{frame}

\begin{frame}{Sphinxを使うきっかけ}
 \begin{itemize}
  \item 4月の宿題の続き。blockdiagシリーズのパッケージング。
	\begin{itemize}
	 \item パッケージング中なう
	 \item 岩松さんにスポンサー依頼
	\end{itemize}
  \item 仕事のプロジェクトでMS Officeの代わりに試験的に導入
 \end{itemize}
\end{frame}

\begin{frame}{Debianでつかう}
 \begin{enumerate}
  \item sudo apt-get install python-sphinx
  \item mkdir hoge;cd hoge; sphinx-quickstart
  \item reSTファイル編集
  \item make html
 \end{enumerate}
\end{frame}

\begin{frame}{日本語環境}
\begin{table}
 \begin{center}
{\scriptsize
  \begin{tabular}{|l|c|}
    \hline
    フォーマット & 結果 \\
    \hline
    html & OK \\
    epub & OK (ただし、CSSは反映されない)\\
    text & OK \\
    man & OK \\
    latex & OK \\
    latexpdf & NG \\
    \hline
  \end{tabular}
}
 \end{center}
\end{table}
生成される \LaTeX がUTF-8…。\\
その他にはGifをインクルードするとコケる。
\end{frame}

\begin{frame}{他の方法でPDFを生成する}
\begin{enumerate}
 \item sudo apt-get install rst2pdf
 \item conf.py編集
 \item Makefile編集
 \item スタイルシート(ja.json)作成
 \item make pdf
\end{enumerate}
\end{frame}

\begin{frame}{問題もあり}
 \begin{itemize}
  \item *diagシリーズのSphinx拡張を使うと、ビルド失敗する
 \end{itemize}
個人的には本末転倒…。
\end{frame}

\begin{frame}{Doxygen}
 \begin{itemize}
  \item ソースコードからドキュメントを生成
  \item C/C++、Java、Python、C\#、Objective-Cなどに対応
  \item 生成可能なフォーマット HTML、 \LaTeX 、RTF(MS-Word)、PostScript、PDF、manなど
 \end{itemize}
\end{frame}

\begin{frame}{Debianで使ってみる}
 \begin{enumerate}
  \item sudo apt-get install doxygen
  \item cd srctree
  \item doxygen -g doxygen.conf
  \item 設定ファイル(doxygen.conf)編集
  \item doxygen doxygen
 \end{enumerate}
\end{frame}

\begin{frame}{Sphinxとの連携}
 \begin{itemize}
  \item breatheを使う\\
	\url{https://github.com/michaeljones/breathe}
 \end{itemize}
\end{frame}

\begin{frame}{まとめ}
 \begin{itemize}
  \item ソースコードからドキュメント生成　: doxygen
  \item reSTからドキュメント生成 : sphinx
  \item doxygen ←→　sphinx : breathe
  \item 宿題がまた増えた
	\begin{itemize}
	 \item rst2pdfで*diagが使えない問題
	 \item breatheのパッケージング
	\end{itemize}
 \end{itemize}
\end{frame}

\end{document}
